% =====================================================================
%  CARNET D'ENTRAÎNEMENT — CHIMIE — FICHE 3 — THÈME 1
%  Particules subatomiques — NIVEAU MOYEN — CORRIGÉ
% =====================================================================
\documentclass[11pt,a4paper]{article}

\usepackage[utf8]{inputenc}
\usepackage[T1]{fontenc}
\usepackage{lmodern}
\usepackage[french]{babel}

\usepackage{amsmath,amssymb,mathtools}
\usepackage{xcolor}
\usepackage{colortbl}
\usepackage{tikz}
\usetikzlibrary{arrows.meta,positioning,calc}
\usepackage[most]{tcolorbox}
\usepackage{enumitem}
\usepackage{booktabs}
\usepackage{array}
\usepackage[a4paper,margin=2.1cm,headheight=26pt,headsep=9pt]{geometry}
\usepackage{fancyhdr}
\usepackage[hidelinks]{hyperref}

\definecolor{primary}{RGB}{124,78,140}
\definecolor{primarydark}{RGB}{74,44,92}
\definecolor{defcol}{RGB}{0,114,178}
\definecolor{thmcol}{RGB}{0,140,120}
\definecolor{excol}{RGB}{0,135,90}
\definecolor{attncol}{RGB}{200,40,55}

\newcommand{\nuc}[3]{\prescript{#1}{#2}{\mathrm{#3}}}
\newcommand{\pp}{p^{+}}
\newcommand{\ee}{e^{-}}
\newcommand{\nz}{n^{0}}
\newcommand{\rep}{\textbf{\color{excol}Réponse : }}

\pagestyle{fancy}
\fancyhf{}
\fancyhead[L]{\small\textcolor{primarydark}{\textbf{Fiche 3 — Thème 1}} : Particules subatomiques}
\fancyhead[R]{\small\textcolor{primarydark}{\textsc{Moyen — Corrigé}}}
\fancyfoot[C]{\small\thepage}
\renewcommand{\headrulewidth}{0.8pt}
\renewcommand{\headrule}{\hbox to\headwidth{\color{primary}\leaders\hrule height \headrulewidth\hfill}}

\tcbset{
  boxbase/.style={enhanced,breakable,boxrule=0.8pt,arc=2pt,
     left=2.6mm,right=2.6mm,top=1.8mm,bottom=1.8mm,
     fonttitle=\bfseries,coltitle=white,toptitle=0.4mm,bottomtitle=0.4mm},
}
\newtcolorbox{bilan}[1][]{boxbase,colback=thmcol!7,colframe=thmcol,
  title={\textsc{À retenir}\ifx\relax#1\relax\relax\else\textmd{ : \itshape #1}\fi}}

\newcounter{exo}
\newcommand{\cor}[1][]{\stepcounter{exo}\par\medskip%
  \noindent{\color{primary}\rule[-2pt]{3.5pt}{15pt}}\ \textbf{\color{primarydark}\large Exercice \theexo}%
  \ \textmd{\normalsize\color{black!55}— corrigé}%
  \ifx\relax#1\relax\relax\else\ \textmd{\normalsize\color{black!65}\itshape (#1)}\fi\par\nopagebreak\smallskip}

\setlist{leftmargin=1.7em,topsep=2pt,itemsep=3pt}

\begin{document}

\begin{tikzpicture}
\node[rectangle,rounded corners=6pt,fill=primarydark,text=white,
      minimum width=\textwidth,minimum height=1.35cm,align=center,inner sep=6pt]
  {\Large\bfseries Thème 1 — Particules subatomiques\\[2pt]
   \normalsize\mdseries Niveau Moyen — Corrigé détaillé};
\end{tikzpicture}

% ---------------------------------------------------------------
\cor[charge d'un grand nombre de protons]
La charge est proportionnelle au nombre de protons :
\[ Q = 10^{12}\times(+1{,}6\times10^{-19}) = (1\times1{,}6)\times10^{\,12+(-19)} = 1{,}6\times10^{-7}\ \text{C}. \]
\rep $\boxed{Q = +1{,}6\times10^{-7}\ \text{C}}$ (charge positive : ce sont des protons).

% ---------------------------------------------------------------
\cor[charge d'un ensemble d'électrons]
Chaque électron porte $-1{,}6\times10^{-19}$ C, donc
\[ Q = 2{,}5\times10^{15}\times(-1{,}6\times10^{-19}) = -(2{,}5\times1{,}6)\times10^{\,15-19}
      = -4{,}0\times10^{-4}\ \text{C}. \]
\rep $\boxed{Q = -4{,}0\times10^{-4}\ \text{C}}$. Le signe est \textbf{négatif} car les électrons portent
une charge négative.

% ---------------------------------------------------------------
\cor[remonter au nombre de particules]
On divise la charge totale par la charge d'un proton :
\[ N = \frac{4{,}8\times10^{-18}}{1{,}6\times10^{-19}}
     = \frac{4{,}8}{1{,}6}\times10^{\,-18-(-19)} = 3\times10^{1} = 30. \]
\rep $\boxed{30 \text{ protons en excès}}$.

% ---------------------------------------------------------------
\cor[rapport des masses]
\begin{enumerate}[label=\alph*)]
  \item $\dfrac{m(\ee)}{m(\pp)} = \dfrac{9{,}1\times10^{-31}}{1{,}7\times10^{-27}}
        = \dfrac{9{,}1}{1{,}7}\times10^{-31-(-27)} \approx 5{,}4\times10^{-4} \approx 0{,}0005$.
        \; \rep l'électron est environ $\mathbf{2000}$ fois plus léger que le proton
        ($1/0{,}0005 = 2000$).
  \item \rep un proton a \textbf{quasiment la même masse} qu'un neutron
        ($\approx 1{,}7\times10^{-27}$ kg tous les deux) ; ni plus lourd, ni plus léger de façon
        significative.
\end{enumerate}

% ---------------------------------------------------------------
\cor[où se cache la masse d'un atome ?]
\begin{enumerate}[label=\alph*)]
  \item Noyau : $12$ nucléons (6 protons + 6 neutrons), chacun $\approx1{,}7\times10^{-27}$ kg :
        \[ m_{\text{noyau}} \approx 12\times1{,}7\times10^{-27} \approx 2{,}0\times10^{-26}\ \text{kg}. \]
  \item Électrons : $6\times9{,}1\times10^{-31} \approx 5{,}5\times10^{-30}\ \text{kg}$.
  \item Fraction : $\dfrac{5{,}5\times10^{-30}}{2{,}0\times10^{-26}} \approx 2{,}7\times10^{-4}
        \approx 0{,}03\,\%$.
        \; \rep les électrons ne pèsent que $\sim3$ dix-millièmes de la masse de l'atome :
        \textbf{la masse est portée quasi exclusivement par le noyau} (les nucléons).
\end{enumerate}

\begin{bilan}
Charge d'un ensemble : $Q = N\times q_{\text{unité}}$ (mantisses $\times$, exposants $+$). Nombre de
particules : $N = Q/q_{\text{unité}}$. La masse d'un atome $\approx$ masse de ses nucléons ; celle des
électrons est négligeable ($\sim0{,}03\,\%$).
\end{bilan}

\end{document}
